% ============================================================================
\section{Method}
\label{sec:method}

\subsection{Retrieval}
For each query the retriever returns a pool of $K=100$ tickets from the allowed
history $H$, and the generator receives the top $k=5$. The generator is called
through the Hugging Face Transformers library \cite{b29}. The primary retriever is
the frozen dense MiniLM bi-encoder \cite{b2,b28} over a FAISS index \cite{b3}; BM25
\cite{b4} and hybrid reciprocal-rank fusion \cite{b5} are used in the matched
comparison of Section~\ref{sec:repair}. All retrievers return pools of the same
size. We do not compare retriever strength under different lift configurations,
because absolute lift units are not comparable across score geometries.

\subsection{Lift}
Let $\mathrm{pos}$ and $\mathrm{neg}$ be the number of positive and negative judge
scores for candidate $c$ in scope $s$, and $n=\mathrm{pos}+\mathrm{neg}$. The
legacy Laplace lift is
\begin{equation}
p=\frac{\mathrm{pos}+\alpha}{n+\alpha+\beta},\qquad
\ell_{\mathrm{L}}=\operatorname{clamp}_{[-c,c]}\!\left(
(p-\tfrac12)\,\min\!\left(1,\tfrac{n}{2}\right)m\right),
\label{eq:legacy}
\end{equation}
with $\alpha=\beta=1$, $m=0.80$ and $c=0.20$. The empirical-Bayes lift is
\begin{equation}
\bar p_s=\pi_s,\quad a=\kappa\pi_s,\quad
p_{\mathrm{post}}=\frac{\mathrm{pos}+a}{n+\kappa},\qquad
\ell_{\mathrm{EB}}=\operatorname{clamp}_{[-c,c]}\!\left(
(p_{\mathrm{post}}-\pi_s)\,m\right),
\label{eq:eb}
\end{equation}
where $\pi_s$ is the mean judge score of scope $s$ over the admissible training
rows, with a global fallback when a scope has no rows. The finalist uses
$\kappa=2$. Two differences from Eq.~\ref{eq:legacy} matter: the neutral point is
the scope base rate rather than $0.5$, and the $n/2$ ramp is absent because the
centring term already attenuates sparse evidence. The legacy ramp is one of the
factors varied in the ablation, not an incidental implementation detail. A
candidate with no evidence receives $\ell=0$.

Pool-relative scaling expresses the bounded lift in retrieval-score units:
\begin{equation}
\ell'=\ell\cdot\frac{\lambda\,\sigma_{\mathrm{pool}}}{c},
\label{eq:pool}
\end{equation}
where $\sigma_{\mathrm{pool}}$ is the standard deviation of the retrieval scores in
the pool and $\lambda=1$ in the pool-scaled cells. The finalist uses absolute
scaling, so Eq.~\ref{eq:pool} is applied only in the ablation. The clamp is applied
before scaling. The parameters $\kappa$ and $\lambda$ were selected on dev over
$\{2,5,10,20,50\}$ and $\{0.25,0.5,1,2,4\}$, respectively; the per-cell values are
in the provenance table. The published legacy cell ($m=0.80$, $c=0.20$) is one
cell in the same sweep, not a special default.

\subsection{Scopes}
The candidate lift is routed by the recorded resolving team and the reply-derived
intent class. Global counts all training queries; Team restricts to the ticket's
team; Class to its intent class; Team$\cap$Class to the intersection, the finest
conventional scope; Backoff uses the finest scope with at least $k_{\min}=3$
observations, falling back from intersection to team to class to Global. The
finalist is Team$\cap$Class. The labels are held fixed across panels so that the
manipulated factor is the data boundary. They are recorded outcomes, not
predictions from the incoming ticket; results are therefore conditional on them,
and Section~\ref{sec:discussion} states this limitation.

\subsection{Control}
The gate is a logistic model over four feature families: intervention shape (how
many candidates the lift would promote, by how much, and the lift dispersion),
baseline quality (top-1 and top-5 scores and margin), semantic evidence
(query--candidate cosine statistics and score--lift correlations), and scope
evidence (coverage and maximum lift per scope). It is trained with grouped
cross-validation by ticket and frozen on dev. We report it as risk
control, not as a source of gain. The hindsight bound applies feedback only when it
helps on the observed outcome and therefore uses information no deployable policy
has; it bounds the headroom, not the performance of a feasible gate. The gate
numbers reported below are from the earlier evaluation on the familiar split; the
strict panel supplies the always-on and hindsight-bound values, not a separately
trained gate. The expected-value analysis is exploratory and is confined to the
Appendix because its rows are two per ticket and are not independent.

\subsection{Evaluation}
The primary outcome is the per-ticket change in cosine between the generated first
reply and the recorded reference reply, averaged over the evaluation tickets. For
ticket $t$,
\begin{equation}
\Delta_t = \cos(\text{feedback reply}_t,\text{ref}_t)
           - \cos(\text{baseline reply}_t,\text{ref}_t),
\label{eq:delta}
\end{equation}
where $\text{ref}_t$ is the recorded reference reply. ROUGE-L \cite{b27} is
reported as a secondary lexical metric. The expert assessment of
Section~\ref{sec:expert} is the independent check of procedural correctness.
Because the study sample contains changed pairs only, we report the conditional
effect and the population effect, the latter as the changed share times the
conditional effect; unchanged pairs have a difference of zero by construction.

The automatic metric compares a generated reply with one recorded reply. It
rewards similarity, not resolution success, and a ticket may admit several valid
replies. BGE cosine \cite{b25} and BERTScore \cite{b26} were computed for the
earlier evaluation; they are not recomputed for every history-only panel because the
offline encoder is slow on CPU and the independent check in this article is the
expert assessment. The Appendix states the difference in coverage. This is a
limitation of the measurement, not a claim that cosine is a complete criterion.
